\section*{Chapter \ref{ch:s1:pairs-and-baskets} --- Pairs and Baskets}

\begin{solution}{exo:s1:pairs-and-baskets:1}
$10 \times \binom{45}{2} = 10 \times 990 = 9\,900$ pairs; a valid 5\% test rejects about $9\,900 \times 0.05 = 495$ of them, all false.
\end{solution}

\begin{solution}{exo:s1:pairs-and-baskets:2}
The procedure rejects the $k$ smallest where the $k$-th is at most $qk/m$: $0.05 \times 20/1000 = 0.001$.
\end{solution}

\begin{solution}{exo:s1:pairs-and-baskets:3}
Opening trades \$1 in each leg and closing trades them again: four units at five basis points, 20 basis points of one leg.
\end{solution}

\begin{solution}{exo:s1:pairs-and-baskets:4}
At $2 \times 2.4\% = 4.8\%$. The rule bets on convergence and has no stop: the takeover premium moved the short leg up by a quarter in a day, and the stock
then stopped trading, so the spread could never come back. The loss is the premium, not the threshold.
\end{solution}

\begin{solution}{exo:s1:pairs-and-baskets:5}
The sum of squared differences is smallest for stocks whose prices move least, whatever their relation. Where substitutes exist they win the ranking (311
of 320 selections); where none exist, the ranking fills with quiet stocks whose paths stayed close by chance, and their spreads drift apart like any
difference of random walks.
\end{solution}

\begin{solution}{exo:s1:pairs-and-baskets:6}
The procedure controls the false discovery rate when the p-values of the true nulls are uniform (or conservative). Here they are not: fat tails and
clustering volatility make the smallest p-values of the 745\,000 null pairs far smaller than a uniform distribution would give, and those false pairs fill the
list of discoveries ahead of the planted ones. The procedure is right; its input is wrong.
\end{solution}

\begin{solution}{exo:s1:pairs-and-baskets:7}
Of the 4\,020 pairs rejected again, 401 are planted, 10.0\%; among the 58\,024 formation-year rejections the planted share is 0.91\%. Requiring persistence
enriches the planted pairs elevenfold: a pair that stays cointegrated is much more likely to be real, which is the named result's use.
\end{solution}

\begin{solution}{exo:s1:pairs-and-baskets:8}
The false-discovery correction assumes calibrated p-values, which the Engle--Granger test does not provide on fat-tailed, heteroskedastic returns; the
survivors are mostly false. And selecting by the same data used to report the Sharpe ratio overstates it. Check the test's calibration on data with no
cointegration, require persistence or an economic link, and report on a later period.
\end{solution}

\begin{solution}{pb:s1:pairs-and-baskets:1}
\begin{enumerate}
\item Long the laggard and short the leader of two stocks that have moved together, closed when they converge; the \omterm{def:s1:pairs-and-baskets:distance}{distance method} selects pairs by the sum of
squared differences of normalised prices over a \omterm{def:s1:pairs-and-baskets:distance}{formation period} and trades them over the next \omterm{def:s1:pairs-and-baskets:distance}{trading period}.
\item Twelve months' formation, six months' trading, open at two formation standard deviations, close at the next crossing, top 20 pairs; up to 11\% a year in
excess returns, 1962--2002.
\item Cointegration: an Engle--Granger test on the log prices, trading the residual; copula: a copula fitted to daily returns, trading an index of improbable
joint moves.
\item 91, 85 and 43 basis points a month before costs for distance, cointegration and copula; 38, 33 and 5 after.
\item Fifty pairs of close substitutes whose log price spread is a stationary AR(1) (ten days, 2\%); without them no pair on the synthetic market is real.
\item Distance and copula 311 of 320, two-step 113, Engle--Granger 40.
\item Twin market: 5.05, 3.89, 1.03, $-0.17$, $-0.49$; default market: $-0.43$, $-0.50$, $-0.11$, $-0.47$, $-0.35$.
\item Planted substitutes selected by distance, opened twice; on day 111 the short leg was taken over, the spread jumped to 35\% and the pair lost 24.6\% per
dollar leg without closing.
\item 745\,465 tests over sixteen cycles; at a nominal 5\% they reject 7.4\% of pairs where none is real, and at 0.1\% three times the nominal rate.
\item 97.2\%: miscalibrated p-values put false pairs ahead of the planted ones.
\item It ranks the 200 closest pairs by p-value, and within them the test still prefers its false rejections: 113 planted pairs, net 1.03.
\item A stock's specific moves have no counterpart in its peers, so the spread against any basket is a random walk.
\item \textbf{Named result.} 7.9\% of the pairs rejected at 5\% in a formation year are rejected again the next year (84.6\% of planted pairs, 7.2\% of the
others); the distance book nets a Sharpe ratio of 5.05 on the market with substitutes and $-0.43$ on the market without.
\item True pairs persist, false ones do so at about the test's false-rejection rate, so persistence separates them.
\item A continuing downward trend in profitability, with strong performance in prolonged turbulence, including the 2008 crisis.
\item Residuals against factor models, structural relative value with contractual links, and broad books of small bets.
\item Run it on data where the null is true (shuffled pairs, or a simulation with the data's tails and volatility) and compare rejection rates with nominal
levels.
\item An economic reason to be a substitute, a spread that reverts on data not used to select it, and liquidity in both legs.
\item Screen for takeover bids and rumours, cap the loss per pair, close pairs whose leg is under offer, and diversify across many pairs.
\item Close substitutes whose prices are pulled back together by something other than their own past.
\end{enumerate}
\end{solution}

\begin{solution}{iq:s1:pairs-and-baskets:1}
Start from an economic link (the same business, share classes, a supplier and customer), then check the link in the data: distance or a cointegration test on
a \omterm{def:s1:pairs-and-baskets:distance}{formation period}, persistence on a later one, liquidity and borrow on both legs. Count every pair tested and hold out time for evaluation.
\end{solution}

\begin{solution}{iq:s1:pairs-and-baskets:2}
Correlation is about returns moving together day by day; cointegration is about prices staying together over time, a stationary spread. Two stocks can be
highly correlated and drift apart, or weakly correlated day to day and still cointegrated. A pairs trade bets on the spread, so cointegration is what matters,
and it is harder to establish.
\end{solution}

\begin{solution}{iq:s1:pairs-and-baskets:3}
Benjamini--Hochberg or a family-wise bound over all tests. What can still go wrong: p-values miscalibrated by fat tails and changing volatility (check them on
null data); dependence between tests; and relationships that are real in the sample and break later. Require out-of-sample persistence.
\end{solution}

\begin{solution}{iq:s1:pairs-and-baskets:4}
The pair's premise is gone: close it, taking the loss if the target is the short leg, unless the book runs merger arbitrage deliberately (chapter~11). Then
check the other pairs for exposure to the same deal or sector.
\end{solution}

\begin{solution}{iq:s1:pairs-and-baskets:5}
With $N$ the $T \times n$ matrix of normalised prices, $D_{ij} = \|N_i\|^2 + \|N_j\|^2 - 2 N_i^\top N_j$: one matrix product $N^\top N$ ($3\,000 \times
3\,000$ from $T \times 3\,000$), then the upper triangle.
\end{solution}

\begin{solution}{iq:s1:pairs-and-baskets:6}
The expected deviation decays as $\phi^h$, so it halves after $h = \ln 2/(-\ln\phi)$ days. The Dickey--Fuller statistic is scale-invariant: its power
depends on $\phi$ and the sample length, not on $\sigma$. A narrow spread is not easier to detect, only cheaper to trade relative to its moves.
\end{solution}
